\answer{ex:22:1}{(a)~$32\%$ et $100\%$. (b)~$R_\uparrow-R_\downarrow\ge2\sqrt{2000\,\text{Hz}/1\,\text{s}}\approx89$~Hz, soit seulement $4.5\%$ de la fréquence totale.}
\answer{ex:22:2}{(a)~$n\ge4(p_1q_1+p_2q_2)/\Delta^2\approx560$. (b)~$\approx56$.}
\answer{ex:22:3}{(a)~Bilan détaillé: $\rho PK$ est symétrique. (b)~La proportion de ratés est la moyenne harmonique $1/\langle1/P\rangle=0.18$, contre $0.5$ sans rétroaction.}
\answer{ex:22:4}{(a)~Un parallélogramme d'aire $4h^2=0.04$ (en tenant compte de la borne sur $p$, numériquement $0.045$). (b)~$\approx0.18$.}
\answer{ex:22:5}{$0.49$: une partie des cultures part dans une région où le raté est presque certain et y erre. Avec la borne: $1.00$; sur un ensemble borné, la densité $\propto1/P$ se concentre dans la région absorbante.}
\answer{ex:22:6}{(a)~Il faut $\ge5$ niveaux; $8$ électrodes suffisent. (b)~Non: il faudrait $r_{\max}\ge25/T=100$~Hz.}
\answer{ex:22:7}{Problème ouvert. Le temps de la recherche aléatoire croît avec $d$ à peu près comme le rapport des volumes $(L/h)^d$, celui de la recherche au signe de l'erreur, linéairement en $d$. Une expérience d'inversion départage les hypothèses: après un raté, un stimulus prévisible mais fort; après un renvoi, un stimulus imprévisible mais faible; il faut quelques dizaines de cultures par groupe.}
